\documentclass{article}
\usepackage{natbib, a4wide, graphicx, amsmath}
\title{Hawkins-Bradu-Kass data}


%\newcommand{\btheta}{\boldsymbol{\theta}}
\newcommand{\epsvec}{\boldsymbol{\epsilon}}
%\newcommand{\deltavec}{\boldsymbol{\delta}}
%\newcommand{\bY}{\mathbf Y}
%\newcommand{\by}{\mathbf y}
%\newcommand{\bx}{\mathbf x}
%\newcommand{\Yvec}{\mathbf{Y}}
%\newcommand{\wvec}{\boldsymbol{\omega}}
%\newcommand{\rvec}{\mathbf r}
%\newcommand{\bone}{\mathbf 1}
%\newcommand{\pW}{p_W}
%\newcommand{\pD}{p_D}
%\newcommand{\pV}{p_V}
%\newcommand{\rE}{\text{E}}


\begin{document}
\maketitle

The Hawkins-Bradu-Kass (HBK) data is an artificial data set
\citep{Hawkins1984} that challenges outlier detection methods based on
perturbing a single observation. Figure~\ref{fig:hbk-data} illustrates
the pairwise relationships between the three predictor variables
$X_1 X_2, X_3$ and outcome $Y$ in the HBK data. The observations can
be clustered into 3 groups: group A, which constitutes the majority of
the observations, shows no relationship between the predictor
variables and the outcome; group B consists of observations that have
high-leverage with respect to group A but maintain the same null
relationship with $Y$; group C are all outliers that pull the fitted
least squares line away from the null.

\begin{figure}
\resizebox{\textwidth}{!}{\includegraphics{hbk-data.pdf}}
\caption{\label{fig:hbk-data} Hawkins-Bradu-Kass (HBK) artificial data
  set. Group A contains the majority of the data; group B is a set of
  4 observations with high leverage; and group C is a set of 10
  outliers. The gray line shows a univariate least-squares fit to each
  predictor variable.}
\end{figure}

\section{Outlier detection}

The masking effect of multiple outliers makes the detection of the
outlying group impossible using univariate methods that consider only
one observation at a time. \citet{Thomas2018} show how the 3 clusters
of observations can be detected after fitting a Bayesian linear model
by applying multivariate analysis methods to the posterior variance
matrix $V$ of the log likelihood contributions.

In a frequentist analysis, outlier-robust linear regression methods
can be used to identify group C as a set of outliers and, after
down-weighting their influence, reveal the underlying null
relationship in groups A and B. In a Bayesian analysis, one way to
improve robustness of a linear model is to use a 2-component normal
mixture for the outcome variable $Y$. This allows a small proportion
of the data to be contaminated with observations with a distinct mean
and variance from the rest. Details of the mixture model are given in
the supplementary materials. After fitting the mixture model,
figure~\ref{fig:hbk-influence} shows the loadings for the first
principal component of $V$ (top panel) and $\Omega$ (bottom
panel). The top panel shows the direction of the perturbation
$\epsvec$ that is maximally influential. This perturbation puts the
highest loads on observations in group C. The bottom panel shows the
maximally outlying perturbation, where observations in group C are
given almost zero weight. The observations in group C are collectively
influential but are not considered outliers in a mixture model that
explicitly accounts for the presence of anomalous observations.

\begin{figure}
\resizebox{\textwidth}{!}{\includegraphics{hbk-influence.pdf}}
\caption{\label{fig:hbk-influence} Diagnostics for a normal mixture
  model fitted to the HBK data. The top panel (influence) shows factor
  loadings for the first principal component of $V$ corresponding to a
  maximally influential perturbation. The bottom panel (outlyingness)
  shows factor loadings for the first principal component of the
  outlier matrix $\Omega$ corresponding to a maximally outlying
  perturbation.}
\end{figure}

\section{Cross conflict}

The HBK data illustrates extreme cross-conflict. If any one of the
observation groups A,B, or C is removed, then the omitted data will be
in contradiction with a linear regression line fitted to the remaining
two groups (See Figure~\ref{fig:hbk-data}).  To quantify this, the
cross-conflict ratio $p_{Vi}/p_{Wi}$ was calculated for each group in
turn. The results are shown in table~\ref{tab:hbk1}. The top row shows
the results for a linear model with a normal error distribution. This
shows very high cross-conflict ratios for all groups with group A
having the highest ratio with $p_{VA}/p_{WA}=59.5$. The second row
shows results for the mixture model 
Cross-conflict ratios are substantially
reduced compared to the linear model but still high for groups A (4.4)
and C (7.9). The use of an outlier-robust mixture model does not
eliminate the cross-conflict, which is a feature of the data, not the
model.

\begin{table}[th]
\centering
\begin{tabular}{lrrr}
\hline
Model          & \multicolumn{3}{c}{Group} \\
\cline{2-4}
               &   A    &   B   &   C  \\
\hline
Normal         &  59.5  &   4.89 & 17.6\\
Normal mixture &   4.4  &   1.81 &  7.9\\
\hline
\end{tabular}
\caption{\label{tab:hbk1} Cross-conflict in the HBK data measured by
  the ratio $p_{Vi}/p_{Wi}$ when only only outcomes from group
  $i \in \left\{A, B, C\right\}$ are considered as data and
  information the other two groups is absorbed into the prior.}
\end{table}

\bibliographystyle{apalike}
\bibliography{HBK}
  
\end{document}
